\documentclass[11pt]{amsart}
\usepackage{amsmath,amssymb}
\setlength{\emergencystretch}{3em}
\tolerance=600
\usepackage{tikz-cd}
\usepackage{tikz}
\usepackage[hidelinks]{hyperref}

\newtheorem{thm}{Theorem}[section]
\newtheorem{lem}[thm]{Lemma}
\newtheorem{cor}[thm]{Corollary}
\newtheorem{prop}[thm]{Proposition}
\theoremstyle{definition}
\newtheorem{defi}[thm]{Definition}
\newtheorem{rem}[thm]{Remark}
\newtheorem{ex}[thm]{Example}
\newtheorem{que}[thm]{Question}

\newcommand{\Coh}{\mathbf{Coh}}
\newcommand{\Mod}{\mathrm{Mod}}
\newcommand{\Set}{\mathbf{Set}}
\newcommand{\Fin}{\mathbf{FinSet}}
\newcommand{\Sub}{\mathrm{Sub}}
\newcommand{\Up}{\mathrm{Up}}
\newcommand{\Ch}{\mathrm{Ch}}
\newcommand{\PP}{\mathsf P}
\newcommand{\HH}{\mathsf H}
\newcommand{\cC}{\mathcal C}
\newcommand{\cD}{\mathcal D}
\newcommand{\LC}{\mathbf{LC}}
\newcommand{\KC}{\mathbf{KC}}

\begin{document}

\title[Intermediate logics and imaginaries]{Intermediate logics and imaginaries: the translation into an intermediate logic commutes with the pretopos completion of all finite propositional theories exactly for the extensions of G\"odel--Dummett logic}
\author{Ryota Shibusawa}
\thanks{Corresponding author. E-mail: \texttt{r-shibusawa@daiichi-koudai.ac.jp}}
\address{Daiichi Institute of Technology}
\email{r-shibusawa@daiichi-koudai.ac.jp}
\keywords{intermediate logic, G\"odel--Dummett logic, weak excluded middle, Heyting category, pretopos completion, imaginaries, Kripke model, G\"odel algebra}
\subjclass[2020]{03G30 (primary), 03B55, 18F10, 06D20 (secondary)}

\begin{abstract}
For an intermediate logic $L$ let $\HH_L\cC$ be the syntactic category
of a coherent theory, with coherent syntactic category $\cC$, read in
first-order logic over $L$, and let $\PP$ be the pretopos completion,
which adds the imaginaries. There is a full and faithful comparison
$\HH_L\PP\cC\to\PP\HH_L\cC$. For intuitionistic logic it is an
equivalence only if $\cC$ is Boolean. We determine how this depends on
$L$. If $L$ does not prove the weak excluded middle $\neg p\vee\neg\neg
p$, the comparison is an equivalence exactly for the Boolean $\cC$. If
$L$ proves it but not the linearity axiom $(p\to q)\vee(q\to p)$, the
comparison fails for the theory of pure equality and for the
propositional theory of the four-element chain. If $L$ proves
linearity, the comparison is an equivalence for every finite
propositional theory. Hence the translation commutes with the
imaginaries of all finite propositional theories if and only if $L$
contains the G\"odel--Dummett logic; for the first-order theory of
pure equality it commutes only for classical logic. The proofs use Kripke models over
the fork and over the diamond, which are frames of $L$ in the first and
second cases, and in the third case the dual of the free $L$-algebra on
a finite distributive lattice, the poset of strict chains of prime
filters of bounded length, over which every finite presheaf embeds in one pulled back
from the lattice.
\end{abstract}

\maketitle

\section{Introduction}

Let $\cC$ be a small coherent category, the syntactic category of a
coherent theory. Its pretopos completion $\PP\cC$ adds finite
coproducts of sorts and quotients by definable equivalence relations,
the imaginaries. In \cite{Heyting} we studied the free Heyting category
$\HH\cC$ on $\cC$, the theory read in intuitionistic first-order logic,
and proved that the canonical comparison $\HH\PP\cC\to\PP\HH\cC$ is full
and faithful and is an equivalence if and only if $\cC$ is Boolean
\cite[Thm.~1.1]{Heyting}. For classical logic, on the other hand, the
corresponding comparison is governed by a Galois criterion and is
often an equivalence \cite{Morley}. The proof for intuitionistic logic
used a Kripke model over the fork, the poset $0<1$, $0<2$, and it was
asked what happens for an intermediate logic, for which the fork need
not be a frame \cite[\S9]{Heyting}. This paper answers the question.

An \emph{intermediate logic} $L$ is a consistent set of propositional
formulas containing the intuitionistic theorems and closed under modus
ponens and substitution. We write $\HH_L\cC$ for the syntactic category
of the internal theory of $\cC$ in intuitionistic first-order logic
with all substitution instances of the formulas of $L$ as axioms
(Section~\ref{sec:HL}), and
\[
c^L_\cC\colon\HH_L\PP\cC\longrightarrow\PP\HH_L\cC
\]
for the canonical comparison, which is full and faithful
(Lemma~\ref{lem:ff}). A \emph{finite propositional theory} is a finite
distributive lattice $D$, regarded as a coherent category: the
lattice of the coherent formulas of a propositional theory, up to
provable equivalence, when this lattice is finite. The \emph{weak excluded middle} is $\neg
p\vee\neg\neg p$ and the \emph{linearity axiom} is $(p\to q)\vee(q\to
p)$; the logics they axiomatise are $\KC$ and the G\"odel--Dummett
logic $\LC\supseteq\KC$.

\begin{thm}\label{thm:main}
Let $L$ be an intermediate logic.
\begin{enumerate}
\item[(1)] If $L$ does not prove the weak excluded middle, then for
every small coherent category $\cC$ the comparison $c^L_\cC$ is an
equivalence if and only if $\cC$ is Boolean.
\item[(2)] If $L$ proves the weak excluded middle and does not prove
the linearity axiom, then $c^L_\cC$ is not an equivalence for the
theory of pure equality and for the four-element chain; but the
objects which obstruct in case \textup{(1)}, two copies of an object
glued along the negation of a subobject, are in the image.
\item[(3)] If $L$ proves the linearity axiom, then $c^L_D$ is an
equivalence for every finite distributive lattice $D$.
\end{enumerate}
In particular $c^L_D$ is an equivalence for all finite distributive
lattices $D$ if and only if $L\supseteq\LC$.
\end{thm}

For first-order theories the situation is different: for the theory of
pure equality $c^L_\cC$ is an equivalence only when $L$ is classical
logic (Theorem~\ref{thm:equality}).

\begin{figure}[h]
\[
\begin{tikzcd}[column sep=small,row sep=small]
1&&2&&&3&\\
&0\arrow[ul]\arrow[ur]&&&1\arrow[ur]&&2\arrow[ul]\\
&&&&&0\arrow[ul]\arrow[ur]&
\end{tikzcd}
\]
\caption{The fork and the diamond. An intermediate logic which does not
contain $\LC$ has one of them as a frame (Lemma~\ref{lem:frames}).}
\end{figure}

\subsection*{What is new and what is known}
New, to the author's knowledge: Theorem~\ref{thm:main}, the embedding
lemma for presheaves over strict chains (Lemma~\ref{lem:kan}), and the
obstruction over the diamond (Proposition~\ref{prop:diamond}), and the
result for pure equality (Theorem~\ref{thm:equality}). Known
and used: the result for intuitionistic logic and its tools
\cite{Heyting}; the duality between finite Heyting algebras and finite
posets with p-morphisms \cite{CZ97}; local finiteness of the variety
of G\"odel algebras \cite{Horn69}; the pretopos completion as the
category of coherent objects of the classifying topos
\cite[D3.3]{Elephant}. Lemma~\ref{lem:frames} is a case of Jankov's
method of characteristic formulas \cite{CZ97}; we give the direct
proof. The description of the free algebra in
Proposition~\ref{prop:free} is known: G\"odel algebras free over
finite distributive lattices are described by chains in
\cite{AGM08}, and free algebras over bounded distributive lattices in
the subvarieties of G\"odel algebras in \cite{Carai}; we include a
proof in the form needed.

\section{The translation into an intermediate logic}\label{sec:HL}

Conventions are those of \cite{Heyting}: $\Coh$ is the
$(2,1)$-category of small coherent categories, a Heyting category is a
coherent category in which the pullback functors of subobjects have
right adjoints $\forall_f$, a Heyting functor preserves them, $\Mod(\cC)$
is the category of models $\cC\to\Set$ and homomorphisms, and universal
properties are used also for targets such as $\Set^K$, the category of
covariant functors from a poset $K$ to sets.

An \emph{$L$-algebra} is a Heyting algebra in which the formulas of $L$
are valid. $L$-algebras are closed under subalgebras, quotients and
products, and every Boolean algebra is an $L$-algebra, $L$ being
contained in classical logic. An \emph{$L$-category} is a Heyting
category $\cD$ such that $\Sub(A)$ is an $L$-algebra for every object
$A$. Boolean coherent categories are $L$-categories.

For $\cC$ in $\Coh$ let $\HH_L\cC$ be the syntactic category of the
internal coherent theory of $\cC$ in intuitionistic first-order logic
(with possibly empty sorts) with the additional axioms
$\lambda(\varphi_1,\dots,\varphi_n)$, for $\lambda(p_1,\dots,p_n)\in L$
and first-order formulas $\varphi_i$ in a common context. It is an
$L$-category, its subobject lattices consisting of formulas in a
context, and the unit $\eta\colon\cC\to\HH_L\cC$ induces equivalences of
groupoids between the Heyting functors $\HH_L\cC\to\cD$ and the coherent
functors $\cC\to\cD$ for every $L$-category $\cD$: a coherent functor
to a Heyting category interprets all first-order formulas soundly for
intuitionistic logic, and the additional axioms are valid in an
$L$-category. Every object of $\HH_L\cC$ is a subobject of some $\eta
A$. For intuitionistic logic, $\HH_L=\HH$.

\begin{lem}\label{lem:basic}
\begin{enumerate}
\item[(a)] If $\cC$ is Boolean, $\eta\colon\cC\to\HH_L\cC$ is an
equivalence.
\item[(b)] If $\cD$ is an $L$-category, $\PP\cD$ is an $L$-category
and a pretopos, the unit $\cD\to\PP\cD$ is Heyting, and a coherent
functor from $\PP\cD$ to a Heyting pretopos is Heyting iff its
restriction to $\cD$ is.
\end{enumerate}
\end{lem}
\begin{proof}
(a) A coherent functor from a Boolean coherent category to a Heyting
category is Heyting \cite[Lemma~2.1]{Heyting}, and $\cC$ is an
$L$-category; so the identity of $\cC$ has the universal property of
$\eta$. (b) Without the words ``$L$-category'' this is
\cite[Lemma~2.2]{Heyting}. An object of $\PP\cD$ is a quotient $X/E$
of a finite coproduct $X=\coprod_iA_i$ of objects of $\cD$, and
$\Sub(X/E)$ is the Heyting subalgebra of the $E$-saturated elements of
$\Sub(X)\cong\prod_i\Sub(A_i)$; so it is an $L$-algebra.
\end{proof}

\begin{lem}\label{lem:ff}
The Heyting functor $c^L_\cC\colon\HH_L\PP\cC\to\PP\HH_L\cC$ extending
$\PP\eta$ is full and faithful, and it is an equivalence if and only if
$\HH_L\PP\cC$ is a pretopos.
\end{lem}
\begin{proof}
As \cite[Lemma~2.3]{Heyting}: by Lemma~\ref{lem:basic}(b) both
$\PP\HH_L\PP\cC$ and $\PP\HH_L\cC$ are the free $L$-category which is a
pretopos on $\cC$, and $c^L_\cC$ is the unit $\HH_L\PP\cC\to\PP\HH_L
\PP\cC$ followed by the resulting equivalence.
\end{proof}

\subsection*{Kripke models over frames of $L$}
A finite poset $K$ is a \emph{frame of $L$} if the Heyting algebra
$\Up(K)$ of its up-sets is an $L$-algebra. An up-set $K'\subseteq K$ of
a frame of $L$ is a frame of $L$, restriction $\Up(K)\to\Up(K')$ being
a surjective homomorphism.

\begin{lem}\label{lem:kripke}
Let $K$ be a poset such that $\uparrow\!k$ is a finite frame of $L$
for every $k\in K$. Then $\Set^K$ is an $L$-category. Hence every
Kripke model $M\colon K\to\Mod(\cC)$ extends to Heyting functors
$M^{\HH_L}$, $M^{\HH_L\PP}$, $M^{\PP\HH_L}$ to $\Set^K$ with
$M^{\PP\HH_L}c^L_\cC\cong M^{\HH_L\PP}$; and for every object $X$ of
$\HH_L\PP\cC$ and all $k\le l$ such that $M(k\le l)$ is an isomorphism
of models, the map $M^{\HH_L\PP}(X)(k)\to M^{\HH_L\PP}(X)(l)$ is
injective.
\end{lem}
\begin{proof}
For a functor $A\colon K\to\Set$ the subobjects of $A$ are the up-sets
of the poset $\int\!A$ of pairs $(k,x)$, $x\in A(k)$, with
$(k,x)\le(l,A(k\le l)x)$; so $\Sub(A)=\Up(\int\!A)$. The restrictions
$\Up(\int\!A)\to\Up(\uparrow\!(k,x))$ are homomorphisms of Heyting
algebras, jointly injective, and $\uparrow\!(k,x)\cong\uparrow\!k$; so
$\Sub(A)$ is a subalgebra of a product of $L$-algebras. The rest is as
in \cite[\S3, Lemma~3.1]{Heyting}: $X$ is a subobject of $\eta Y$ with
$Y$ in $\PP\cC$, and the transition maps of $M^{\PP}(Y)$ along
isomorphisms of models are bijections.
\end{proof}

\section{Which frames an intermediate logic has}

\begin{lem}\label{lem:frames}
Let $L$ be an intermediate logic.
\begin{enumerate}
\item[(a)] The fork is a frame of $L$ if and only if $L$ does not
prove the weak excluded middle.
\item[(b)] If $L$ proves the weak excluded middle and not the linearity
axiom, the diamond is a frame of $L$.
\item[(c)] If $L$ proves the linearity axiom, a finite poset $Z$ is a
frame of $L$ if and only if every $\uparrow\!z$ is a chain with at most
$n_L$ elements, where $n_L\in\{1,2,\dots,\infty\}$ is the supremum of
the cardinalities of the finite chains which are frames of $L$.
\end{enumerate}
\end{lem}
\begin{proof}
Let $B$ be the Lindenbaum algebra of $L$, an $L$-algebra in which a
formula is valid iff it is in $L$.

(a) In $\Up(\text{fork})$ the weak excluded middle fails for
$p=\{1\}$, since $\neg p=\{2\}$ and $\neg\neg p=\{1\}$. Conversely let
$b\in B$ with $\neg b\vee\neg\neg b\ne1$ and $c=\neg b$, so that
$\neg\neg c=c$. The set $\{0,c,\neg c,c\vee\neg c,1\}$ is a Heyting
subalgebra of $B$: $c\wedge\neg c=0$, $c\to\neg c=\neg c$, $\neg c\to
c=\neg\neg c=c$, $(c\vee\neg c)\to c=c$, $(c\vee\neg c)\to\neg c=\neg
c$ and $\neg(c\vee\neg c)=0$. Its five elements are distinct because
$c\vee\neg c\ne1$, and it is isomorphic to $\Up(\text{fork})$, with $c$
and $\neg c$ the two atoms. So $\Up(\text{fork})$ is an $L$-algebra.

(b) Let $a,b\in B$ with $(a\to b)\vee(b\to a)\ne1$, and $x=a\to b$,
$y=b\to a$, $t=x\wedge y$. In every Heyting algebra $x\to y=y$ and
$y\to x=x$: if $z\wedge(a\to b)\le b\to a$ then $z\wedge b\le a$, as
$b\le a\to b$. Hence $\{t,x,y,x\vee y,1\}$ is closed under $\wedge$,
$\vee$ and $\to$: $(x\vee y)\to x=y\to x=x$, $(x\vee y)\to y=y$,
$x\to t=x\to y=y$, $y\to t=x$ and $(x\vee y)\to t=x\wedge y=t$. Let $B'=B/F$ be the quotient by the
filter $F$ generated by $\neg\neg t$. In $B'$ the class of $\neg t$ is
$0$, since $\neg t\wedge\neg\neg t=0$; so the classes of
$0,t,x,y,x\vee y,1$ form a Heyting subalgebra $G$ of $B'$, the
negation of each of $t,x,y,x\vee y$ being $0$ in $B'$. We show that
the class of $x\vee y$ is not $1$, that is, $\neg\neg t\not\le x\vee
y$. In $B$ the weak excluded middle gives $\neg t\vee\neg\neg t=1$ and
the De Morgan law $\neg(x\wedge y)=\neg x\vee\neg y$ (for $\neg(x\wedge
y)\wedge\neg\neg x\le\neg y$ in every Heyting algebra, and $\neg
x\vee\neg\neg x=1$). Now $\neg x=\neg(a\to b)=\neg\neg a\wedge\neg
b\le\neg b\le y$ and likewise $\neg y\le x$; so $\neg t\le x\vee y$,
and $\neg\neg t\le x\vee y$ would give $x\vee y=1$. It follows that
the six classes are distinct: the classes of $x$ and $y$ are
incomparable, since $[x]\le[y]$ gives $[y]=[x\to y]=1$; so $[t]$ is
strictly below both and $[x\vee y]$ strictly above; and $[t]\ne0$,
since $[t]=0$ iff $t\wedge\neg\neg t=0$, while $t\wedge\neg\neg t=t$ and $t=0$ would give $\neg\neg t=0\le
x\vee y$. The algebra $G$ is isomorphic to $\Up(\text{diamond})$, with
$[t]=\{3\}$, $[x]=\{1,3\}$, $[y]=\{2,3\}$. So $\Up(\text{diamond})$ is
an $L$-algebra.

(c) Let $Z$ be a finite frame of $L$ and $z\in Z$. Then $\uparrow\!z$
is a frame of $L$; if it had incomparable elements $a,b$, the up-sets
$p=\uparrow\!a$ and $q=\uparrow\!b$ would give $z\notin(p\to
q)\cup(q\to p)$. So $\uparrow\!z$ is a chain, with at most $n_L$
elements. Conversely a chain with $m\le n_L$ elements is an up-set of
a chain which is a frame of $L$, hence a frame of $L$; and $\Up(Z)$ is
a subalgebra of the product of the $\Up(\uparrow\!z)$.
\end{proof}

\section{Without weak excluded middle: only Boolean categories}

\begin{proof}[Proof of Theorem~\ref{thm:main}\textup{(1)}]
If $\cC$ is Boolean, so is $\PP\cC$, and $\HH_L\PP\cC\simeq\PP\cC$ is a
pretopos by Lemma~\ref{lem:basic}(a); so $c^L_\cC$ is an equivalence
by Lemma~\ref{lem:ff}. Let $S\le A$ be a subobject of $\cC$ which is
not complemented. By \cite[Lemma~4.1]{Heyting} there are a
homomorphism of models $h\colon M\to N$ and $a\in A(M)\setminus S(M)$
with $h_A(a)\in S(N)$. Let $Q$ be the object of $\PP\HH_L\cC$ obtained
from two copies of $A$ by gluing along the pseudo-complement $\neg S$,
and $M^\ast$ the Kripke model over the fork with $M^\ast_0=M^\ast_1=M$,
$M^\ast_2=N$, the identity along $0<1$ and $h$ along $0<2$. By
Lemma~\ref{lem:frames}(a) the fork and its up-sets are frames of $L$,
so Lemma~\ref{lem:kripke} applies. As in the proof of
\cite[Thm.~1.1]{Heyting}, $a\in(\neg M^\ast(S))(1)$ and
$a\notin(\neg M^\ast(S))(0)$, so the two copies of $a$ are distinct
points of $M^{\ast\PP\HH_L}(Q)(0)$ with the same image at $1$, which is
impossible for an object in the image of $c^L_\cC$.
\end{proof}

\section{With weak excluded middle but without linearity}

When $L$ proves the weak excluded middle, the object $Q$ of the
preceding proof is in the image of $c^L_\cC$: in
$\Sub_{\HH_L\cC}(\eta A)$ one has
$\neg S\vee\neg\neg S=A$ with $\neg S\wedge\neg\neg S=0$, so $A$ is the
coproduct of $\neg S$ and $\neg\neg S$ in $\PP\HH_L\cC$, and
\[
Q\;\cong\;\neg S+\neg\neg S+\neg\neg S,
\]
a subobject of $A+A+A$ defined by formulas, hence the image of an
object of $\HH_L\PP\cC$. The obstruction moves from the negation to
the implication.

\begin{prop}\label{prop:diamond}
Let $L$ be an intermediate logic of which the diamond is a frame, and
$\cC$ a small coherent category. Suppose that there are subobjects
$S,T\le A$ in $\cC$, homomorphisms of models $M\xrightarrow{h}N
\xrightarrow{h'}N'$ and $a\in A(M)$ such that
\[
a\notin S(M),\qquad h(a)\in S(N)\setminus T(N),\qquad h'h(a)\in T(N').
\]
Then $c^L_\cC$ is not an equivalence: the object of $\PP\HH_L\cC$
obtained from two copies of $A$ by gluing along $S\Rightarrow T$, the
Heyting implication in $\Sub_{\HH_L\cC}(\eta A)$, is not
in its essential image.
\end{prop}
\begin{proof}
Let $U=(\eta S\Rightarrow\eta T)\le\eta A$ in $\HH_L\cC$ and $Q=(A+A)/E$
in $\PP\HH_L\cC$, with $E$ the equivalence relation on $A+A$ which is
the union of the diagonal and of the two images of $U$ in the
off-diagonal summands $A\times A$ of $(A+A)\times(A+A)$, as in
\cite[\S5]{Heyting}. Let
$M^\ast$ be the Kripke model over the diamond with
\[
M^\ast_0=M^\ast_1=M,\qquad M^\ast_2=N,\qquad M^\ast_3=N',
\]
the identity along $0<1$, $h$ along $0<2$, $h'$ along $2<3$ and $h'h$
along $1<3$. The up-sets of the diamond are frames of $L$, so
Lemma~\ref{lem:kripke} applies, and $M^{\ast\PP\HH_L}(Q)(k)$ is the
quotient of two copies of $A(M^\ast_k)$ identifying the copies of the
points of $U_{M^\ast}(k)$, where, the implication of $\Set^K$ being the
Kripke implication,
\[
U_{M^\ast}(k)=\{x\mid\text{for all }l\ge k,\ x|_l\in S(M^\ast_l)
\text{ implies }x|_l\in T(M^\ast_l)\}.
\]
For $k=1$ and $x=a$ the nodes $l\ge1$ are $1$, where $a\notin S(M)$,
and $3$, where $h'h(a)\in T(N')$; so $a\in U_{M^\ast}(1)$. For $k=0$
the node $l=2$ has $h(a)\in S(N)\setminus T(N)$; so $a\notin
U_{M^\ast}(0)$. Hence the two copies of $a$ are distinct at $0$ and
equal at $1$, along a transition which is an isomorphism of models,
and $Q$ is not in the essential image of $c^L_\cC$ by
Lemma~\ref{lem:kripke}.
\end{proof}

\begin{proof}[Proof of Theorem~\ref{thm:main}\textup{(2)}]
By Lemma~\ref{lem:frames}(b) the diamond is a frame of $L$. For the
theory of pure equality, with one sort $\sigma$, take
$A=\sigma^3$, $S=[x_1=x_2]$, $T=[x_1=x_2\wedge x_2=x_3]$, a
three-element set $M$ with $a$ an enumeration of it, and maps $M\to
N\to N'$ onto a two-element set, identifying $a_1$ and $a_2$, and then
onto a one-element set. For the four-element chain $0<q<p<1$, a
propositional theory, take $A=1$, $S=p$, $T=q$ and the prime filters
$\{1\}\subset\{p,1\}\subset\{q,p,1\}$. Proposition~\ref{prop:diamond}
applies in both cases. The statement on the objects of case (1) was
shown before the proposition.
\end{proof}

\section{With linearity: finite propositional theories}

Throughout this section $L$ proves the linearity axiom and $n=n_L$ is
as in Lemma~\ref{lem:frames}(c).

\subsection*{Finite duality}
For a finite poset $X$, $\Up(X)$ is a finite Heyting algebra, and every
finite Heyting algebra is of this form, with $X$ the poset of its
prime filters. A monotone map $f\colon Z\to X$ induces the lattice
homomorphism $f^{-1}\colon\Up(X)\to\Up(Z)$, and every homomorphism of
bounded lattices $\Up(X)\to\Up(Z)$ is of this form for a unique $f$;
$f^{-1}$ is a homomorphism of Heyting algebras if and only if $f$ is a
\emph{p-morphism}, $f(\uparrow\!z)=\uparrow\!f(z)$ for all $z$
\cite{CZ97}. A finite distributive lattice $D$ is $\Up(X_D)$ for the
poset $X_D$ of its prime filters, which are the models of $D$ as a
coherent category, ordered by inclusion, the homomorphisms.

For a finite poset $X$ let $\Ch_n(X)$ be the set of the nonempty
strict chains $c=(c_0<c_1<\dots<c_m)$ of $X$ with at most $n$ elements,
ordered by $c\le c'$ iff $c'$ is a final segment of $c$, and
$\pi\colon\Ch_n(X)\to X$, $\pi(c)=c_0$, a monotone map. The up-set of
$c$ is the chain of its final segments.

\begin{prop}\label{prop:free}
For a finite distributive lattice $D=\Up(X)$, the free $L$-algebra on
the bounded lattice $D$ is $\Up(\Ch_n(X))$ with $\pi^{-1}\colon D\to
\Up(\Ch_n(X))$; and $\HH_LD$ is equivalent to it, as a preorder.
\end{prop}
\begin{proof}
$\Up(\Ch_n(X))$ is an $L$-algebra by Lemma~\ref{lem:frames}(c). Let
$\Up(Z)$ be a finite $L$-algebra and $f\colon Z\to X$ monotone. For
$z\in Z$ the up-set $\uparrow\!z$ is a chain $z=z_0<\dots<z_r$ with at
most $n$ elements, and $f(z_0)\le\dots\le f(z_r)$; let $g(z)$ be the
strict chain of the distinct values. Then $g\colon Z\to\Ch_n(X)$ is a
p-morphism with $\pi g=f$: the final segments of $g(z)$ are the
$g(z_i)$. It is the only one: if $g'$ is a p-morphism with $\pi g'=f$,
the elements of the chain $g'(z)$ are the $\pi(c')$, $c'\ge g'(z)$,
that is the $\pi g'(z')=f(z')$, $z'\ge z$. So $\Up(\Ch_n(X))$ has the
universal property for finite $L$-algebras; applied to the subalgebra
generated by $\pi^{-1}(D)$ and to itself, it shows that $\pi^{-1}(D)$
generates it. An arbitrary $L$-algebra is a G\"odel algebra, and
G\"odel algebras are locally finite \cite{Horn69}; so a lattice
homomorphism from $D$ to an $L$-algebra has its image in a finite
subalgebra, and the universal property holds for all $L$-algebras.
Finally every object of $\HH_LD$ is a subobject of some $\eta d$, a
subterminal object; so $\HH_LD$ is a preorder, its poset of subterminal
objects is an $L$-algebra, an $L$-algebra is an $L$-category, and
Heyting functors between Heyting algebras are homomorphisms; the
universal properties agree.
\end{proof}

\subsection*{Pretopos completions}
For a finite poset $X$ the category $\Fin^X$ of functors $X\to\Fin$ is
a pretopos, the functor $\Up(X)\to\Fin^X$ sending an up-set to its
characteristic subterminal functor is coherent, and it is the pretopos
completion of $\Up(X)$: $\Set^X$ is the topos of sheaves on $X$ with
the topology of up-sets, the classifying topos of the propositional
theory $\Up(X)$, and its coherent objects are the finite functors
\cite[D3.3]{Elephant}. For $f\colon Z\to X$ the functor $\PP(f^{-1})$
is $f^*\colon\Fin^X\to\Fin^Z$, composition with $f$. Every $Y$ in
$\Fin^X$ is covered by subterminal objects: $y\in Y(x)$ defines a
morphism $s_y\colon\uparrow\!x\to Y$, and the $s_y$ are jointly
surjective.

So, with $Z=\Ch_n(X)$,
\[
\PP D=\Fin^X,\qquad\PP\HH_LD=\Fin^Z,\qquad\PP\eta=\pi^*,
\]
and $\Fin^Z$ is an $L$-category by Lemma~\ref{lem:kripke}.

\begin{lem}[embedding]\label{lem:kan}
Every $W$ in $\Fin^Z$, $Z=\Ch_n(X)$, is a subobject of $\pi^*Y$ for
some $Y$ in $\Fin^X$; namely for the left Kan extension $Y$ of $W$
along $\pi$, the unit $W\to\pi^*Y$ is a monomorphism.
\end{lem}
\begin{proof}
$Y(x)$ is the colimit of $W$ over the poset $(\pi\!\downarrow\!x)$ of
the chains $c$ with $c_0\le x$. For such $c$ the elements $c_i\le x$
form an initial segment of $c$; let $\tau_x(c)$ be the final segment
of $c$ beginning at the last of them. The chains $c'\ge c$ in
$(\pi\!\downarrow\!x)$ are the final segments of $c$ beginning at some
$c_i\le x$; so $\tau_x(c)$ is the greatest of them and
$\tau_x(c')=\tau_x(c)$. Hence each connected component of
$(\pi\!\downarrow\!x)$ has a greatest element, and
\[
Y(x)=\coprod_{r\in R_x}W(r),\qquad R_x=\{c\mid c_0\le x,\text{ and }
c=(c_0)\text{ or }c_1\not\le x\},
\]
a finite set, the map $W(c)\to Y(x)$ being $W(c\le\tau_x(c))$ into the
summand of $\tau_x(c)$. For $x=\pi(c)=c_0$ the chain $c$ is in $R_x$,
because $c=(c_0)$ or $c_1>c_0$; so the unit $W(c)\to Y(\pi(c))$ is a coproduct
injection.
\end{proof}

For a logic which does not prove linearity the up-sets of the points of
the dual of the free algebra need not be chains, and the existence of
these greatest elements is not guaranteed; this is where the proof
uses linearity.

\begin{proof}[Proof of Theorem~\ref{thm:main}\textup{(3)}]
By Lemma~\ref{lem:ff} it suffices to show that $c=c^L_D\colon\HH_L\PP
D\to\PP\HH_LD=\Fin^Z$ is essentially surjective. We have $c\,\eta_{\PP
D}\cong\pi^*$, and $c\circ\HH_L(\iota)\cong\iota_{\HH_LD}$, where
$\iota$ denotes the units of $\PP$, both sides being Heyting functors
$\HH_LD\to\Fin^Z$ extending the same coherent functor on $D$.

Let $W$ be in $\Fin^Z$. By Lemma~\ref{lem:kan}, $W$ is a subobject of
$\pi^*Y=c(\eta Y)$ for some $Y$ in $\Fin^X=\PP D$. Let $s_j\colon
u_j\to Y$, $j=1,\dots,m$, be morphisms of $\PP D$ from subterminal
objects $u_j\in D$ which are jointly surjective. For each $j$ the
pullback $\varphi_j:=(\pi^*s_j)^*W$ is a subobject of $\pi^*u_j$ in
$\Fin^Z$, a subterminal object, that is an element of $\Up(Z)=\HH_LD$
below $\pi^{-1}(u_j)$; let $\varphi'_j\le\eta u_j$ be its image under
$\HH_L(\iota)\colon\HH_LD\to\HH_L\PP D$, and
\[
X:=\bigvee_j\exists_{\eta s_j}\varphi'_j\ \le\ \eta Y
\]
in $\HH_L\PP D$. The functor $c$ preserves images and unions and sends
$\varphi'_j$ to $\varphi_j$, so $c(X)=\bigvee_j\exists_{\pi^*s_j}
(\pi^*s_j)^*W=W$, the $\pi^*s_j$ being jointly surjective.
\end{proof}

\begin{proof}[Proof of the last assertion of Theorem~\ref{thm:main}]
If $L\supseteq\LC$, part (3) applies. If not, $L$ does not prove
linearity; by parts (1) and (2), $c^L_D$ is not an equivalence for the
four-element chain $D$, which is not Boolean.
\end{proof}

\section{First-order theories: pure equality only for classical logic}

The positive result of the preceding section does not extend to
first-order theories. Already for the theory without symbols and
axioms the translation commutes with the imaginaries only when the
logic is classical. The obstruction is the quotient of the sort by the
double negation of equality, and the frame is the two-element chain.

\begin{lem}\label{lem:two}
The two-element chain $0<1$ is a frame of every intermediate logic
other than classical logic.
\end{lem}
\begin{proof}
If $L$ is not classical, its Lindenbaum algebra has an element $b$ with
$d:=b\vee\neg b\ne1$. Then $\neg d=\neg b\wedge\neg\neg b=0$, so
$\{0,d,1\}$ is a Heyting subalgebra, isomorphic to the algebra of the
up-sets of the two-element chain.
\end{proof}

\begin{thm}\label{thm:equality}
Let $\cC$ be the syntactic category of the coherent theory of pure
equality, with one sort $\sigma$ and no symbols and axioms. For an
intermediate logic $L$ the comparison $c^L_\cC$ is an equivalence if
and only if $L$ is classical logic.
\end{thm}
\begin{proof}
For classical logic $\HH_L$ is the Morleyization and the statement is
\cite[Ex.~6.5]{Morley}. Let $L$ be non-classical, and $K$ the chain
$0<1$, a frame of $L$ with its up-sets (Lemma~\ref{lem:two}). The
models of $\cC$ are the sets and the homomorphisms are all maps. Let
$\Omega$ be a countably infinite set with a partition into at least
two blocks, all of them infinite, $p\colon\Omega\to\Omega/{\sim}$ the projection to
the set of blocks, and $M$ the Kripke model over $K$ with $M_0=\Omega$,
$M_1=\Omega/{\sim}$ and transition $p$.

In every Heyting category $\neg\neg(x=y)$ is an equivalence relation on
$\sigma$; let $W=\sigma/\neg\neg(x=y)$ in $\PP\HH_L\cC$. In $\Set^K$ a
pair $(a,b)$ of elements of $\Omega$ satisfies $\neg\neg(x=y)$ at $0$
iff $p(a)=p(b)$, by the formula for the Kripke negation; so
$M^{\PP\HH_L}(W)(0)$ is the set of blocks.

Suppose that $W\cong c^L_\cC(X)$ with $X\le\eta Y$ in $\HH_L\PP\cC$,
$Y$ in $\PP\cC$. Write $Y=(\coprod_iA_i)/F$ with objects $A_i$ of
$\cC$, definable subsets of powers $\sigma^{n_i}$ in the language of
equality, and $F$ a coherent equivalence relation. By
Lemma~\ref{lem:kripke} we obtain an injection
\begin{multline*}
\theta\colon\{\text{blocks}\}=M^{\PP\HH_L}(W)(0)\ \cong\ M^{\HH_L\PP}
(X)(0)\\ \subseteq\ Y(\Omega)=\Big(\coprod_iA_i(\Omega)\Big)/F(\Omega).
\end{multline*}
Let $G$ be the group of the permutations of $\Omega$ which preserve the
partition. Each $g\in G$, with the induced permutation of the blocks,
is an automorphism of the Kripke model $M$; it induces natural
automorphisms of the extensions of $M$ (an extension together with its
identification with $M$ on $\cC$ being unique up to a unique
isomorphism), and the isomorphism $W\cong c^L_\cC(X)$ is natural
with respect to them. So $\theta$ is $G$-equivariant, $G$ acting on
the blocks and on $Y(\Omega)$ in the evident way. Moreover every
permutation $\pi$ of $\Omega$ is an automorphism of the model $\Omega$
and preserves $F(\Omega)$.

Let $B$ be a block and $\theta(B)$ the class of a tuple $d\in
A_i(\Omega)$. Since the blocks are infinite there is $g\in G$ mapping
every block to itself with $g(d_j)\notin\{d_1,\dots,d_{n_i}\}$ for all
$j$; then $gB=B$, so $F(d,gd)$. Let $B'\ne B$ be another block; there
is $\varepsilon\in G$ exchanging $B$ and $B'$, mapping the other
blocks to themselves, with $\varepsilon(d_j)\notin\{d_1,\dots,
d_{n_i}\}$ for all $j$. The map $d_j\mapsto d_j$, $g(d_j)\mapsto
\varepsilon(d_j)$ is a well defined injective map between finite
subsets of $\Omega$, so it extends to a permutation $\pi$ of $\Omega$,
and $F(d,gd)$ gives $F(\pi d,\pi gd)=F(d,\varepsilon d)$. Hence
$\varepsilon\theta(B)=\theta(B)$, that is $\theta(\varepsilon B)=
\theta(B)$ with $\varepsilon B=B'\ne B$, contradicting the injectivity
of $\theta$. (If $n_i=0$ the class $\theta(B)$ is fixed by $G$, with
the same contradiction.)
\end{proof}

The transition $p$ of this Kripke model is not an isomorphism, and the
invariant of Lemma~\ref{lem:kripke} is not what excludes $W$; the
argument is the one which shows in \cite[Thm.~5.3]{Morley} that a
classical imaginary is not coherent, an automorphism which moves a
class and cannot be told from one which fixes it by a coherent
formula. For propositional theories there are no such automorphisms,
which is consistent with Theorem~\ref{thm:main}(3).

\section{Summary and questions}

\begin{center}
\small\begin{tabular}{l|l|l}
$L$ proves & frame used & $c^L_\cC$ is an equivalence\\\hline
no w.e.m. & fork & iff $\cC$ is Boolean\\
w.e.m., no linearity & diamond & not for pure equality, the 4-chain\\
linearity & chains ($\le n_L$) & for finite distributive lattices\\
not classical & two-element chain & not for pure equality\\
classical logic & point & iff the criterion of \cite[Thm.~6.1]{Morley}
\end{tabular}
\end{center}

Here w.e.m.\ is the weak excluded middle. The last line is included for comparison: for classical logic
$\HH_L$ is the Morleyization, every finite propositional theory
satisfies the criterion \cite[Prop.~5.2]{Morley}, and the theory of
pure equality satisfies it as well \cite[Ex.~6.5]{Morley}.

\begin{que}
Let $L\supseteq\LC$ be different from classical logic. By
Theorem~\ref{thm:equality} the comparison fails for pure equality. For
which coherent categories $\cC$ is $c^L_\cC$ an equivalence? For a finite
distributive lattice $D$ the pretopos $\cC=\PP D$ is an example which is
not a preorder, by Theorem~\ref{thm:main}(3); but it is Morita
equivalent to the propositional theory $D$. Are there examples which
are neither Boolean nor Morita equivalent to a propositional theory,
that is, with $\PP\cC$ not equivalent to $\PP D$ for any distributive
lattice $D$?
\end{que}

\begin{que}
Does Theorem~\ref{thm:main}\textup{(3)} hold for all bounded
distributive lattices? Both sides are compatible with directed unions
of sublattices if $\HH_L$ and $\PP$ are; we have not checked the
details.
\end{que}

\begin{que}
In case \textup{(2)}, for which non-Boolean $\cC$ is $c^L_\cC$ an
equivalence? The three-element chain has no configuration as in
Proposition~\ref{prop:diamond}.
\end{que}

\section*{Declarations}
The author has no competing interests to declare.

\begin{thebibliography}{99}
\bibitem{AGM08} S.~Aguzzoli, B.~Gerla, V.~Marra, \emph{G\"odel algebras
free over finite distributive lattices}, Ann. Pure Appl. Logic 155
(2008), 183--193.
\bibitem{Carai} L.~Carai, \emph{Free algebras and coproducts in
varieties of G\"odel algebras}, preprint, arXiv:2406.05480.
\bibitem{CZ97} A.~Chagrov, M.~Zakharyaschev, \emph{Modal Logic}, Oxford
Logic Guides 35, Oxford Univ. Press, 1997.
\bibitem{Elephant} P.~T.~Johnstone, \emph{Sketches of an Elephant: A Topos
Theory Compendium}, Oxford Univ. Press, 2002.
\bibitem{Horn69} A.~Horn, \emph{Logic with truth values in a linearly
ordered Heyting algebra}, J. Symbolic Logic 34 (1969), 395--408.
\bibitem{Heyting} R.~Shibusawa, \emph{The free Heyting category on a
pretopos is a pretopos only in the Boolean case: the intuitionistic
translation and the pretopos completion}, preprint (2026), archived at
DOI 10.5281/zenodo.24148558.
\bibitem{Morley} R.~Shibusawa, \emph{Translations in the lattice of
two-out-of-three classes: Morleyization as a sameness of theories, its
join with Morita equivalence, and the forced order of translation and
completion}, preprint (2026), archived at DOI 10.5281/zenodo.24122457.
\end{thebibliography}

\end{document}
